\documentclass[11pt]{article}
\usepackage[a4paper,margin=2.5cm]{geometry}
\usepackage[T1]{fontenc}
\usepackage[utf8]{inputenc}
\usepackage{amsmath,amssymb}
\usepackage{booktabs,array,longtable}
\usepackage{xcolor}
\usepackage{listings}
\usepackage[numbers,sort&compress]{natbib}
\usepackage[colorlinks=true,linkcolor=blue!50!black,citecolor=blue!50!black,urlcolor=blue!50!black]{hyperref}
\lstset{basicstyle=\ttfamily\small,columns=fullflexible,frame=single,framerule=0.3pt,
  rulecolor=\color{black!30},xleftmargin=4pt,xrightmargin=4pt}

\newcommand{\q}[1]{qg_{#1}}
\newcommand{\qZ}{qg_Z}
\newcommand{\qX}{qg_X}
\newcommand{\qY}{qg_Y}

\title{Quantum gates and algorithms in qg units:\\
Hamming-weight conservation and readout classes}
\author{Vicente Humberto Monteverde\\
\small Universidad del Museo Social Argentino (UMSA), Argentina\\
\small ORCID 0000-0001-8884-4811}
\date{\small Draft of 1 October 2026}

\begin{document}
\sloppy
\maketitle

\begin{abstract}
We write fifteen standard quantum gates and fourteen standard quantum algorithms in qg
units~\cite{monteverde2026qang}. An $n$-qubit state is described by its qg values,
$\q{P}=\langle P\rangle$ for every Pauli string $P$, and a gate $U$ acts as $\q{P}'=\q{U^\dagger PU}$. With this one
rule each gate becomes a table of qg transformations and each algorithm a readout identity: which qg values
hold the answer. The formulation is the Pauli (Heisenberg) representation and adds no new physics. It
gives three practical classifications in the units used by the rest of the qang project: which gates
conserve Hamming weight, and therefore admit the qg symmetry witness and filter; which algorithms are read
out from a few qg values and which need the full outcome histogram; and where reading qg values is strictly
worse than the standard readout. Three readout identities are, as far as we know, not stated in this
form elsewhere: the per-qubit polar bias after Grover iterations, Simon's secret as the unique non-trivial
Z-parity equal to one, and quantum counting as a single $\qX$. Every identity is checked numerically and
pinned by tests in the open-source library \texttt{qang} (version 0.5.1, \texttt{pip install qang}),
module \texttt{qang.formulation}.
\end{abstract}

\section{Why write gates and algorithms in qg units}

The qang project describes a qubit by measurement-level quantities, the polar bias $\qZ=\langle\sigma_z\rangle$
and its companions~\cite{monteverde2026qang}, and builds diagnostics on them: a register-mean witness and a
shot-level filter for circuits that conserve a Hamming weight~\cite{monteverde2026witness}, and decoders that
read relaxation from the final readout~\cite{monteverde2026leakage}. Those tools apply to some circuits and not
to others. The question of this note is the bookkeeping one: for the gates and algorithms that people actually
use, where do the qg tools apply, and what does a qg readout give? We answer it with the Pauli representation,
which is standard~\cite{gottesman1998,nielsen2010quantum}, stated in qg units and checked numerically line by
line.

\section{Framework}

\paragraph{State.} For one qubit, $\rho=\tfrac12(I+\qX X+\qY Y+\qZ Z)$; for $n$ qubits,
\begin{equation}\label{eq:state}
\rho=\frac{1}{2^n}\sum_{P\in\{I,X,Y,Z\}^{\otimes n}}\q{P}\,P,\qquad \q{P}=\mathrm{Tr}(\rho P),
\end{equation}
and the purity is $\mathrm{Tr}\rho^2=2^{-n}\sum_P\q{P}^2$. Strings with one non-identity letter give the local qg
values; the others are correlations, such as $\q{ZZ}$.

\paragraph{Measurement.} A Z-basis measurement gives, from the same shots, every diagonal value
$\q{Z_S}=\sum_y p(y)(-1)^{S\cdot y}$ for every subset $S$ of qubits; for one qubit $P(0)=(1+\qZ)/2$. Values with X or
Y need a basis rotation first.

\paragraph{Gates.} A gate $U$ maps qg values as
\begin{equation}\label{eq:rule}
\q{P}'=\q{U^\dagger PU}.
\end{equation}
For a Clifford gate, $U^\dagger PU$ is a signed Pauli string and the gate permutes qg values with signs; for a
non-Clifford gate it is a combination, and the gate mixes them. Convention: the leftmost letter of a string is
qubit 0, which is the control of CX and CZ (and the controls of Toffoli and Fredkin).

\paragraph{Weight conservation.} If $U$ maps each Hamming-weight sector into itself, the register mean
$\overline{\qZ}=1-2N/n$ stays fixed, and the qg witness and filter apply.

\section{Fifteen gates}

\begin{longtable}{>{\raggedright\arraybackslash}p{2.4cm}>{\raggedright\arraybackslash}p{6.6cm}>{\raggedright\arraybackslash}p{5.0cm}}
\caption{Action of each gate on qg values, Eq.~\eqref{eq:rule}; $c=\cos\theta$, $s=\sin\theta$. Values not listed
are unchanged.}\label{tab:gates}\\
\toprule
gate & action on qg values & reading \\
\midrule
\endfirsthead
\toprule
gate & action on qg values & reading \\
\midrule
\endhead
1. Pauli X & $(\qX,\qY,\qZ)\to(\qX,-\qY,-\qZ)$ & flips the polar bias \\
2. Pauli Y & $(\qX,\qY,\qZ)\to(-\qX,\qY,-\qZ)$ & flips $\qZ$ and $\qX$ \\
3. Pauli Z & $(\qX,\qY,\qZ)\to(-\qX,-\qY,\qZ)$ & invisible in the Z basis \\
4. Hadamard & $(\qX,\qY,\qZ)\to(\qZ,-\qY,\qX)$ & swaps polar bias and $\qX$ \\
5. Phase S & $(\qX,\qY)\to(-\qY,\qX)$ & quarter turn on the equator \\
\phantom{5.} $P(\varphi)$ & $\qX'=\cos\varphi\,\qX-\sin\varphi\,\qY$, $\qY'=\sin\varphi\,\qX+\cos\varphi\,\qY$ & same as $R_z(\varphi)$ \\
6. T ($\pi/8$) & $\qX'=(\qX-\qY)/\sqrt2$, $\qY'=(\qX+\qY)/\sqrt2$ & first gate that mixes: irrational qg values (``magic'') \\
7. $R_x(\theta)$ & $\qZ'=c\,\qZ+s\,\qY$, $\qY'=c\,\qY-s\,\qZ$ & rotation in the YZ plane \\
8. $R_y(\theta)$ & $\qZ'=c\,\qZ-s\,\qX$, $\qX'=c\,\qX+s\,\qZ$ & from $|0\rangle$: $\qZ=\cos\theta$, the qang identity; the RQang gate is $R_y(\arccos\qZ)$ \\
9. $R_z(\theta)$ & $\qX'=c\,\qX-s\,\qY$, $\qY'=s\,\qX+c\,\qY$ & invisible in the Z basis \\
10. CNOT & $\q{XI}'=\q{XX}$, $\q{IZ}'=\q{ZZ}$, $\q{IY}'=\q{ZY}$, $\q{YI}'=\q{YX}$ & turns a correlation into a local value: the target's $\qZ$ becomes the parity $\q{ZZ}$ (syndrome extraction) \\
11. CZ & $\q{XI}'=\q{XZ}$, $\q{IX}'=\q{ZX}$, $\q{YI}'=\q{YZ}$, $\q{IY}'=\q{ZY}$ & invisible in the Z basis; conserves weight \\
12. SWAP & $\q{ab}'=\q{ba}$ & relabels; conserves weight \\
13. iSWAP & $\q{ZI}'=\q{IZ}$, $\q{XI}'=-\q{ZY}$, $\q{YI}'=\q{ZX}$, $\q{IX}'=-\q{YZ}$, $\q{IY}'=\q{XZ}$ & swap with phase; conserves weight \\
14. Toffoli & $\q{IIZ}'=\tfrac12(\q{IIZ}+\q{IZZ}+\q{ZIZ}-\q{ZZZ})$ & mixes inside the diagonal sector \\
15. Fredkin & $\q{IZI}'=\tfrac12(\q{IZI}+\q{ZZI}+\q{IIZ}-\q{ZIZ})$ & controlled swap; conserves weight \\
\bottomrule
\end{longtable}

Three observations organize the table.
\begin{itemize}
\item \textbf{Clifford or not.} X, Y, Z, H, S, CNOT, CZ, SWAP and iSWAP permute qg values with signs. T, generic
rotations, Toffoli and Fredkin mix them: one string opens into two (T) or four (Toffoli, Fredkin) per
application. That growth is what makes non-Clifford circuits hard to simulate~\cite{gottesman1998}.
\item \textbf{Z-basis blindness.} Z, S, T, $R_z$, $P(\varphi)$ and CZ leave every $\qZ$ fixed; their effect appears
only after a gate that moves the equator to the poles. This is, gate by gate, the Z-basis blind spot of the qg
metrics~\cite{monteverde2026qang}.
\item \textbf{Weight conservation.} Among multi-qubit gates, CZ, SWAP, iSWAP and Fredkin conserve Hamming weight;
CNOT and Toffoli do not. A circuit built only from weight-conserving gates and $R_z$ rotations has zero sector
exposure, and in our simulations the qg filter then removes almost all relaxation error; basis changes around
two-qubit gates let about half of it through~\cite{monteverde2026witness}.
\end{itemize}

\section{Fourteen algorithms}

For each: the readout identity, what reading it in qg gives, and what it does not.

\paragraph{1. Deutsch--Jozsa~\cite{deutsch1992}.} With $f$ as a phase oracle, $P(0^n)=2^{-n}\sum_S\q{Z_S}$: the
all-zero probability is the mean of all Z-parities. In qg: $f$ is constant if and only if all $n$ local $\qZ$
equal $+1$; one value below 1 already means balanced. No correlation is needed.

\paragraph{2. Bernstein--Vazirani~\cite{bernstein1997}.} The output is the product state $|s\rangle$, so
$\qZ^{(i)}=(-1)^{s_i}$: each secret bit is the sign of its qubit's polar bias. Under noise each bit is decided
separately, with its own interval.

\paragraph{3. Simon~\cite{simon1997}.} The output is uniform over $\{y:y\cdot s=0\}$, so $\q{Z_S}=1$ for
$S\in\{0,s\}$ and $0$ otherwise: the secret is the unique non-trivial Z-parity equal to one. This describes the
answer but does not find it efficiently, since there are $2^n$ parities; the efficient route remains a linear
system from $O(n)$ samples.

\paragraph{4. Shor~\cite{shor1997}.} The quantum part estimates the phase of an order $r$. If $r$ is a power of
two, the local values show it: for $r=4$ with six counting qubits, $\qZ=(0,0,1,1,1,1)$. For $r=6$ the local $\qZ$
lie between 0 and 0.33 and do not determine $r$; the information is in the correlations, and the full bit string
plus continued fractions is needed. This is a negative result for local readout.

\paragraph{5. Grover~\cite{grover1996}.} With one marked item $x^\ast$ among $N=2^n$ and success probability $P$
after $k$ iterations,
\begin{equation}\label{eq:grover}
\qZ^{(i)}=(-1)^{x^\ast_i}\,\frac{NP-1}{N-1}.
\end{equation}
The marked item is the sign pattern of the local polar biases, and their common magnitude measures success.
Reading it this way costs more shots than the histogram, about $1/P^2$ against $1/P$, as we found in a noisy
simulation; the identity is a diagnostic, not a better readout.

\paragraph{6. Phase estimation~\cite{kitaev1995} and 11. phase kickback.} A control in $|+\rangle$ acting on an
eigenstate of $U$ with eigenvalue $e^{i\varphi}$ ends with
\begin{equation}\label{eq:kick}
\qX+i\,\qY=e^{i\varphi},\qquad \qZ\ \text{unchanged}.
\end{equation}
Kickback writes the phase into the control's equatorial angle. Iterative phase estimation reads $\qX$ and $\qY$
for $2^k\varphi$, one bit at a time; it is the natural qg form of phase estimation, and the same mechanism drives
Deutsch--Jozsa and Bernstein--Vazirani.

\paragraph{7. HHL~\cite{harrow2009}.} The output $|x\rangle\propto A^{-1}|b\rangle$ cannot be read out in full; what
HHL delivers is expectation values $\langle x|M|x\rangle$, that is qg values of the solution register, and its
success probability is the ancilla's polar bias, $P(1)=(1-\qZ^{\mathrm{anc}})/2$. The qg form states plainly
what HHL produces: qg values, not the vector~\cite{aaronson2015}. We check this readout for a $2\times2$ matrix by
exact linear algebra, not with a circuit.

\paragraph{8. Quantum Fourier transform.} The QFT of a basis state $|x\rangle$ is a product of equatorial
qubits: qubit $l$ has $\qZ=0$ and $\qX+i\,\qY=e^{2\pi i x/2^l}$. The QFT moves the information from polar biases
(the bits of $x$) to equatorial angles, which the Z basis does not see; hence it is always followed by a basis
change before measurement.

\paragraph{9. VQE~\cite{peruzzo2014}.} The energy is linear in qg values, $E=\sum_P c_P\,\q{P}$, grouped by
measurement basis. This is the main use of the qg tools so far: witness, filter, symmetry checks and pole-damped
optimization~\cite{monteverde2026qang,monteverde2026witness}.

\paragraph{10. QAOA~\cite{farhi2014}.} For MaxCut, $C=\sum_{(i,j)}w_{ij}(1-\q{Z_iZ_j})/2$: the whole cost comes from
ZZ correlations in the same shots. With a weight-conserving XY mixer the witness and filter
apply~\cite{monteverde2026witness}.

\paragraph{12. Quantum counting~\cite{brassard1998}.} A Hadamard test of the Grover operator $G$ on the uniform
state $|s\rangle$ leaves the control with
\begin{equation}\label{eq:count}
\qX=\langle s|G|s\rangle=1-\frac{2M}{N}\quad\Longrightarrow\quad M=\frac{N(1-\qX)}{2}.
\end{equation}
Counting solutions is measuring one $\qX$: with $N=16$ and three marked items, $\qX=0.625$ and $M=3$. The precision
follows single-qubit statistics; full counting by phase estimation gains precision with more counting qubits.

\paragraph{13. Quantum walks~\cite{farhi1998walk}.} A continuous-time single walker with hopping $(XX+YY)/2$
conserves the number of excitations, so the register mean stays at $\overline{\qZ}=1-2/n$ at all times (checked
for $n=5$): single-walker walks live in the weight-one sector, where the witness and filter apply as in Heisenberg
dynamics. Discrete coined walks do not conserve weight.

\paragraph{14. Machine learning: Q-SVM~\cite{havlicek2019}, Q-PCA~\cite{lloyd2014qpca}.} The fidelity kernel is an
all-zero probability, $k(x,x')=P(0^n)=2^{-n}\sum_S\q{Z_S}$, and for product angle encodings it reduces to
$\prod_i(1+\vec q_i\cdot\vec q_i\,')/2$, the qang transition probability. In Q-PCA the purity and spectrum of $\rho$
depend on squared qg values, $\mathrm{Tr}\rho^2=2^{-n}\sum_P\q{P}^2$.

\section{Readout classes}

\begin{table}[h]
\centering\small
\caption{Which qg values suffice to read out each algorithm.}
\label{tab:readout}
\begin{tabular}{>{\raggedright\arraybackslash}p{4.4cm}>{\raggedright\arraybackslash}p{9.8cm}}
\toprule
sufficient readout & algorithms \\
\midrule
local $\qZ$ & Deutsch--Jozsa, Bernstein--Vazirani \\
one ancilla's $\qX$, $\qY$ & phase kickback, iterative phase estimation, quantum counting (Hadamard test), HHL success \\
few-qubit correlations & QAOA (ZZ), VQE (Hamiltonian strings) \\
all parities or the histogram & Simon, Shor with general order, Grover (local readout exists, at $\sim1/P^2$ shots), fidelity kernels \\
\bottomrule
\end{tabular}
\end{table}

\section{What is new and what is not}

The Pauli representation of states, the action of gates by conjugation, and the stabilizer formalism are
known~\cite{gottesman1998,nielsen2010quantum}, and so are the algorithms~\cite{deutsch1992,bernstein1997,simon1997,
shor1997,grover1996,kitaev1995,harrow2009,peruzzo2014,farhi2014,brassard1998,farhi1998walk,havlicek2019,
lloyd2014qpca}. The contribution is organizational and practical: one language, the qg units of the qang
project, for three questions. (i) Which gates conserve weight and admit the qg witness and filter. (ii) Which
algorithms are read out from a few qg values and which need the histogram. (iii) Where qg readout is worse
(Grover, Shor with general order, Simon). The readout forms in Eqs.~\eqref{eq:grover} and~\eqref{eq:count}, and
Simon's secret as the unique unit parity, are elementary but, to our knowledge, not stated in this form
elsewhere. None of this is a quantum advantage.

\section{Limitations}

(i) Ideal, noiseless states throughout; noise enters only through the cited qg studies. (ii) The identities are
checked on small instances ($n\le6$). (iii) The HHL readout is checked by exact linear algebra for a $2\times2$
matrix, not with a circuit. (iv) Shot costs are quoted from a separate simulation (Grover) or are the standard
single-qubit statistics.

\section{Code availability}
The library \texttt{qang} (MIT license) is on PyPI (\texttt{pip install qang}, version 0.5.1) and at
\url{https://github.com/Viny2030/qang}. The module \texttt{qang.formulation} gives the gate matrices, the rule of
Eq.~\eqref{eq:rule} (\texttt{conjugate}, \texttt{gate\_table}, \texttt{apply\_gate}), qg values of a state
(\texttt{qg\_values}) and one readout function per algorithm; \texttt{tests/test\_formulation.py} pins every
identity, and \texttt{notebooks/qang\_inicio\_formulacion.ipynb} runs them in Colab. For example:
\begin{lstlisting}
from qang import formulation as F
bell = F.apply_gate(F.apply_gate(F.qg_values([1, 0, 0, 0]),
                                 F.kron(F.H, F.I2)), F.CX)
# qg_XX = 1, qg_YY = -1, qg_ZZ = 1, all local values 0
\end{lstlisting}

\section*{Acknowledgments and use of AI tools}
Generative AI tools (Anthropic Claude Opus 5.5) assisted with code, numerical checks and drafting. The author
framed the questions, made the design decisions, and reviewed and validated all AI-assisted output.

\bibliographystyle{unsrtnat}
\bibliography{../refs}
\end{document}
